% Ported from thesis/05-results.md — wording unchanged.
\chapter{Results}\label{ch:results}

% Numbers in this chapter are FINAL (post-#245): the 7-session campaign
% (s01–s07) is the primary H1–H3 evidence; the pooled pre-campaign run-set is
% pilot-only. Every number traces to an archived run via Appendix A.
% Wording per scope-of-claims.md.

Reading order: H1 establishes that the aggregate instrument is blind; H2--H4
show placement reproducibly moves mechanism layers that do not reach the
user; H5--H6 locate where placement does bite --- node failure --- as a measured
latency$\leftrightarrow$availability trade-off. The fault-class grouping is churn (H1--H3),
load (H4--H5), node failure (H6).

\begin{figure}[tb]
  \centering
  \includegraphics[width=\textwidth]{fig-02-core-matrix.png}
  \caption{The fault-class $\times$ measurement-layer core matrix: which layer
    moves under churn / load / node failure, and whether it reaches the user.}
  \label{fig:core-matrix}
\end{figure}

\section{H1 --- The aggregate score cannot rank placement strategies}\label{sec:h1}

\textbf{Statement.} The probe-based aggregate resilience score does not
reproducibly discriminate placement strategies; between-strategy differences
are a small fraction of total score variance and undetectable at feasible
iteration counts.

\textbf{Result --- supported (7-session campaign, 147 churn iterations).}

\begin{center}
\begin{tabular}{p{0.5\textwidth}p{0.42\textwidth}}
\toprule
\textbf{Quantity} & \textbf{Value} \\
\midrule
\texttt{ICC\_strategy} & \textbf{0.033} (cluster-bootstrap 95\% CI \textbf{[0.014, 0.178]}) \\
Variance partition & between-strategy \textbf{3.3\%} / run-to-run \textbf{37.6\%} / iteration \textbf{59.1\%} \\
Focal contrast (colocate 64.0 vs spread 74.3) & $d = 0.46$ $\rightarrow$ \textbf{73 iterations/strategy} for 80\% power \\
MDE at the $n = 3$ actually run per session & 2.29 sd $\approx$ \textbf{51 score points} --- larger than any gap that exists \\
\bottomrule
\end{tabular}
\end{center}

The earlier pooled run-set (16 mixed-version runs: ICC 0.046, focal $d = 0.06$
needing $\approx$3,982/strategy) read the same way and is retained as the pilot only.

Scope (per \S\ref{sec:positioning}): this is ``cannot \textbf{rank placement strategies under session
variance}'' --- MicroRes's 0.86--0.90 \emph{binary} classification accuracy is a
different task and is not contradicted.

\begin{figure}[tb]
  \centering
  \includegraphics[width=\textwidth]{fig-03-h1-score-distributions.png}
  \caption{Per-strategy score distributions across the seven sessions
    (overlapping spreads, means clustered 64--74).}
  \label{fig:h1-dist}
\end{figure}

\begin{figure}[tb]
  \centering
  \includegraphics[width=\textwidth]{fig-04-h1-icc-trajectory.png}
  \caption{ICC trajectory as sessions accumulate, 0.222 $\rightarrow$ 0.033
    (run-to-run variance accrues with k).}
  \label{fig:h1-icc}
\end{figure}

\section{H2 --- Placement reproducibly moves a kernel/network reconvergence signature}\label{sec:h2}

\textbf{Statement.} Under churn, spreading the target's dependents across nodes
flushes a large fraction of per-node conntrack state during the kill cycle;
co-location does not.

\textbf{Result --- supported (7/7 sessions).} Median conntrack flush: \texttt{spread}
\textbf{38.5\%} vs \texttt{colocate} \textbf{2.7\%}; \texttt{spread > colocate} in \textbf{7/7} independent
sessions --- \textbf{sign test $p = 0.0156$, paired Wilcoxon $W = 0$, $p = 0.0225$}.
The pooled pilot agreed (36.6\% vs 1.9\%, 16/16 runs). This is the most
reproducible signal in the study.

Secondary, corroborating only: \texttt{colocate} throttles CPU lowest in 6/7
sessions (the pooled pilot had \texttt{best-fit} lower still) --- lead with conntrack.

\textbf{Mechanism --- protocol composition measured.} A dedicated probe (per-node
\texttt{conntrack -L} protocol counts at 5 s intervals through one full kill cycle
under each placement; \S\ref{sec:h2-mechanism}, raw data + chaos-window timestamps in
\texttt{data/conntrack-probe/}) yields three window-robust composition findings:
(1) \textbf{TCP dominates the table under both placements and drops sharply at
the kill cycles in both} (spread 5,935 $\rightarrow$ 4,253, $-$28\%; colocate
4,973 $\rightarrow$ 3,937, $-$21\%) --- kernel-side teardown of flows traversing the killed
pod: kube-proxy has no TCP conntrack-deletion path (\#100698, \#104098;
verified across the $\le$v1.31 exec cleaner and $\ge$v1.32 netlink reconciler), and
close-state conntrack timeouts (10--120 s; \S\ref{sec:h2-mechanism}) expire torn-down flows
within the cycle;
(2) \textbf{the clearly placement-dependent component is UDP (DNS)}: under
steady load \texttt{spread} sustains $\sim$\textbf{4$\times$} more UDP entries than \texttt{colocate}
(chaos-window medians 910 vs 224), consistent with cross-node calls driving
connection churn and DNS re-resolution, and with kube-proxy's documented
UDP-only cleanup (\#48370, \#108523, \#126130; ipvs mode here) having more to
clean under spread; (3) with $i = 1$ and a Locust-ramp-contaminated
baseline, the probe \textbf{cannot apportion the campaign's flush percentages}
between the two paths --- both mechanisms are visible; their shares remain
unquantified. The seven sessions carry the statistics; the probe carries
composition and event timing.

\begin{figure}[tb]
  \centering
  \includegraphics[width=\textwidth]{fig-05-h2-conntrack.png}
  \caption{Conntrack flush \% per strategy across the seven sessions
    (spread > colocate in all seven; note \texttt{default} flushes as hard as \texttt{spread}
    ($\sim$42\%) while \texttt{dependency-aware}/\texttt{best-fit} show \emph{negative} flush --- entries
    grow during chaos --- consistent with the flush tracking placement-dependent
    state composition rather than a fixed per-strategy quantity).}
  \label{fig:h2-conntrack}
\end{figure}

\section{H3 --- The mechanism is decoupled from the user-visible outcome}\label{sec:h3}

\textbf{Statement.} The reproducible mechanism (H2) does not translate into a
reproducible user-visible outcome on the fault-dependent routes beyond a
run-level confound.

\textbf{Result --- decoupling supported (campaign + three pilot tests).} Across the
7 sessions (49 strategy-cells): conntrack flush $\rightarrow$ \textbf{dependent}-route p95 is
\textbf{$\rho = 0.07$ ($p = 0.65$)} and \textbf{decoupled by TOST}, while the \textbf{control}
route shows $\rho = 0.29$ ($p = 0.043$) --- the mechanism correlates with the route
that does \emph{not} depend on the killed service: the signature of a run-level
confound, not propagation. The only dependent-significant association
(TCP-retransmit delta, $\rho = -0.32$) is \emph{negative} --- opposite a propagation
story. Pilot agreement: pooled $\rho = 0.15$ n.s.; the one significant mechanism
(CoreDNS p99) stronger on control (0.54) than dependent (0.31); within-run
mean $\rho \approx +0.10$; robust to route reclassification.

No statistics needed for the headline table: \texttt{dependency-aware} has the
\emph{worst} mechanism (conntrack +20\%) and the \emph{best} dependent-route error rate
(1.4\%); \texttt{spread} flushes 9$\times$ more than \texttt{colocate} yet they tie on user-visible
error (8.0\% vs 8.9\%).

Why would a kernel-layer signature this reproducible fail to reach the user
routes? The plausible mechanism reading: under single-replica churn the
user-visible outcome is dominated by the outage itself --- the target's only
replica is gone for the same deletion-to-ready cycle whatever the topology ---
while client retries and timeouts mask the differences in reconvergence
timing that the conntrack signature represents. This interpretation is ours
(per the convention of Chapter~\ref{ch:discussion}); what is measured is the layer split, and
\S\ref{sec:decoupling} develops the account in full.

\begin{figure}[tb]
  \centering
  \includegraphics[width=\textwidth]{fig-06-h3-scatter.png}
  \caption{Mechanism-vs-outcome scatter (dependent vs control routes
    per strategy-cell): the flush correlates with the route that does \emph{not}
    depend on the killed service.}
  \label{fig:h3-scatter}
\end{figure}

\section{H4 --- Under load contention, placement moves the mechanism, not (reproducibly) the user}\label{sec:h4}

This section contains the first of two results the provenance discipline
retracted --- the demonstration that the gate works as designed (\S\ref{sec:campaign}).

\textbf{Result --- mechanism replicates; user layer does not.} Two $i = 4$ batches
(the second with fully clean, doctor-gated provenance):

\begin{itemize}
\item \emph{Replicated}: east-west inter-service p95, median spread/colocate ratio
  \textbf{1.39$\times$ (batch A)} and \textbf{1.36$\times$ (batch B)} --- co-location keeps calls
  node-local.
\item \emph{Not replicated}: user-facing during-load ratios collapsed from $\sim$2.1--2.4$\times$
  (batch A, dependent > control) to $\sim$1.05--1.40$\times$ (batch B, dependent $\approx$
  control --- no dependency specificity). The dirty pilot's ``co-location $\sim$3$\times$
  better at the user layer'' reading did not survive replication.
\end{itemize}

\textbf{No user-visible placement effect is claimed under load.} The finding
matches the churn result: layered decoupling holds across both fault classes.

Why does the east-west effect replicate while the user routes do not? The
inter-service hop is the direct mechanical path placement changes --- a call
either crosses a node boundary or it does not --- whereas a user-facing route
aggregates many hops and, under a 200-user spike, sits behind cluster-wide
queueing noise that affects every route alike. The locality term plausibly
drowns in that first-order saturation before it reaches the front end. This
interpretation is ours (per the convention of Chapter~\ref{ch:discussion}); \S\ref{sec:tradeoff} prices the
east-west locality face that does replicate against its availability
counterpart.

\section{H5 --- A graph-derived metric separates node-local from spreading placements}\label{sec:h5}

This section contains the second result the provenance discipline retracted:
replication kept the separation but took back the headline correlation (\S\ref{sec:campaign}).

\textbf{Result --- the two-regime separation replicates across two independent
batches; a continuous law does not.} Across all 8 strategies under a
200-user spike ($i = 4$ per batch):

\begin{center}
\begin{tabular}{lrrrr}
\toprule
\textbf{strategy} & \textbf{frac (batch 1)} & \textbf{EW p95 (batch 1)} & \textbf{frac (batch 2)} & \textbf{EW p95 (batch 2)} \\
\midrule
\textbf{colocate} & 0.00 & \textbf{33.9} & 0.00 & \textbf{33.2} \\
\textbf{best-fit} & 0.13 & \textbf{35.3} & 0.00 & \textbf{35.7} \\
dependency-aware & 0.73 & 42.6 & 0.73 & 43.8 \\
spread & 0.73 & 43.5 & 0.73 & 44.2 \\
baseline & 0.70 & 43.5 & 0.78 & 41.7 \\
adversarial & 0.80 & 43.5 & 0.80 & 42.0 \\
default & 0.78 & 45.5 & 0.82 & 41.6 \\
random & 0.80 & 43.9 & 0.80 & 42.8 \\
\bottomrule
\end{tabular}
\end{center}

In \textbf{both} batches the two node-local placements (cross-node fraction $\approx$ 0)
occupy the two lowest east-west tails of eight --- per-batch null probability
$1/28 \approx 0.036$, jointly $\approx 0.0013$ --- at $\sim$\textbf{1.25$\times$} below the spreading cluster
($\sim$34.6 vs $\sim$43.5 ms medians in batch 1; $\sim$34.5 vs $\sim$42.4 ms in batch 2). The
fraction, computed pre-chaos from the dependency graph + actual
\texttt{podPlacements}, captures that separation.

The \emph{continuous} correlation did \textbf{not} replicate: batch 1's Spearman
$\rho = 0.79$ ($n = 8$, carried by \texttt{best-fit}'s intermediate 0.13 point) collapsed
to \textbf{$\rho = 0.25$ (n.s.)} in batch 2, where \texttt{best-fit} packed fully (fraction
0.00) and the spreading cluster showed no internal trend. We therefore claim
H5 as a \textbf{replicated two-regime separator}, not a smooth predictor --- and
quote $\rho = 0.79$ only alongside $\rho = 0.25$.

Secondary findings: locality is \textbf{not unique to \texttt{colocate}} (\texttt{best-fit}'s
bin-packing lands node-local in both batches --- any node-packing placement
gets the benefit); \textbf{\texttt{dependency-aware} did not deliver} (fraction 0.73 in
both batches --- the BFS partition did not co-locate communicating services as
intended; its tail $\approx$ spread's both times).

Caveats: the user layer stays weak ($\sim$1.3$\times$, not dependency-specific); batch
2's launching tree was dirty in non-code files only (deck binary + figure
PNGs; running code = \texttt{e543fbb}); the aggregate score is \textbf{saturated} in
these load runs (batch 1: all eight strategies scored 100; CIs overlap in
both batches --- H1 again, and the saturation \S\ref{sec:decoupling} discusses); this is
\textbf{validation of a static separator}, not of the locality concept (\S\ref{sec:positioning}).

\begin{figure}[tb]
  \centering
  \includegraphics[width=\textwidth]{fig-07-h5-fraction-vs-tail.png}
  \caption{Cross-node fraction vs during-load east-west p95 (batch 1;
    square markers = the node-local pair). Batch 2 reproduces the separation with
    \texttt{best-fit} at fraction 0.00; see the table above.}
  \label{fig:h5-fraction}
\end{figure}

\section{H6 --- Co-location is a latency/availability trade-off}\label{sec:h6}

\textbf{Result --- supported and reproduced (two doctor-clean node-drain batches:
$i = 1$ and $i = 3$).} Draining the node hosting \texttt{productcatalogservice}:

\begin{center}
\begin{tabular}{llll}
\toprule
\textbf{placement} & \textbf{services on drained node} & \textbf{blast radius (observed)} & \textbf{target recovery (mean)} \\
\midrule
\textbf{colocate} & 11 / 11 & \textbf{11 --- whole app offline} & \textbf{$\approx$ 10.3 s} \\
\textbf{spread} & 2 / 11 & \textbf{2 (18\%)} & \textbf{$\approx$ 2.6 s} \\
\bottomrule
\end{tabular}
\end{center}

Observed blast equals placement-predicted blast in every iteration, measured
from EndpointSlice outage troughs (15 s sampling) --- not the score, which is
unusable here (a drain leaves every Litmus probe \texttt{Unknown}; H1 again).

\textbf{Gradient extension (6 placing strategies $\times$ $i = 3$, doctor-clean).}
The full gradient run confirms the relationship is exact across the
intermediate placements: \textbf{observed blast equals the placement-predicted
blast for every strategy} ---

\begin{center}
\begin{tabular}{llll}
\toprule
\textbf{strategy} & \textbf{services on drained node} & \textbf{blast (observed)} & \textbf{recovery (mean)} \\
\midrule
colocate & 11 & \textbf{11} & 10.8 s \\
random & 4 & \textbf{4} & 4.6 s \\
dependency-aware & 3 & \textbf{3} & 33.3 s \\
best-fit & 3 & \textbf{3} & 9.7 s \\
spread & 2 & \textbf{2} & 2.6 s \\
adversarial & 2 & \textbf{2} & 33.1 s \\
\bottomrule
\end{tabular}
\end{center}

Spearman(predicted, observed) = \textbf{1.0} ($n = 6$): per-node concentration
predicts the availability consequence exactly, the availability analogue of
H5's separator. \textbf{Recovery time, however, is \emph{not} monotone in blast} ---
intermediate-blast placements produced both the fastest (4.6 s) and slowest
(33.3 s) recoveries --- so the recovery claim is scoped to the
colocate-vs-spread extremes contrast (10.3 s vs 2.6 s, $\sim$4$\times$), where it
reproduced across both original batches. (Recovery times are not pooled
across batches: the gradient run is a distinct batch and reads colocate at
10.8 s vs 10.3 s in the two-point batches.)

\textbf{The trade-off is the finding}: the same co-location that gives the lowest
east-west tail (H5: $\sim$33--34 ms in both batches) gives a 100\% single-drain
outage; \texttt{spread} is the mirror. One graph property, two opposing measured
consequences. Framing per scope-of-claims: this is the \textbf{quantification of a
known qualitative trade-off}, not its discovery (the full novelty framing
lives in \S\ref{sec:tradeoff}) --- the empirical content is that it materializes under real
chaos with no partial survival, holds across the full strategy gradient, and
drives a measured recovery penalty at the extremes.

Caveats: single-replica, single cluster; recovery non-monotonicity above.

\begin{figure}[tb]
  \centering
  \includegraphics[width=\textwidth]{fig-08-h6-trough-timeline.png}
  \caption{EndpointSlice ready-count trajectory through the drain
    (colocate vs spread; the $i = 3$ batch catches the trough that a fixed-offset
    snapshot missed --- the motivation for the 15 s trough sampler).}
  \label{fig:h6-trough}
\end{figure}

\begin{figure}[tb]
  \centering
  \includegraphics[width=\textwidth]{fig-09-tradeoff.png}
  \caption{The headline figure: the same placements plotted on the
    east-west latency axis (H5) and the node-drain blast-radius axis (H6),
    opposing gradients of co-location.}
  \label{fig:tradeoff}
\end{figure}

\section{Per-claim evidence table}\label{sec:evidence-table}

\begin{center}
\footnotesize
\begin{tabular}{p{0.24\textwidth}p{0.15\textwidth}p{0.21\textwidth}p{0.07\textwidth}p{0.21\textwidth}}
\toprule
\textbf{Claim} & \textbf{Data} & \textbf{Test} & \textbf{Figure} & \textbf{Archived run(s) (Appendix~\ref{app:provenance})} \\
\midrule
H1: score cannot rank under session variance & s01--s07, 147 iterations & variance partition; ICC + cluster-bootstrap CI; power/MDE & \ref{fig:h1-dist}, \ref{fig:h1-icc} & run-20260608-233543 \ldots{} run-20260610-130249 (7 archives) \\
\addlinespace
H2: conntrack flush, spread > colocate & s01--s07 & sign test (7/7, $p = 0.0156$); Wilcoxon ($W = 0$, $p = 0.0225$) & \ref{fig:h2-conntrack} & same 7 archives \\
\addlinespace
H3: mechanism $\perp$ dependent-route outcome & s01--s07, 49 cells & Spearman dep 0.07 vs ctrl 0.29*; TOST equivalence & \ref{fig:h3-scatter} & same 7 archives \\
\addlinespace
H4: east-west replicates, user layer does not & 2 $\times$ $i = 4$ load batches & ratio replication (1.36--1.39$\times$); dep-vs-ctrl collapse & --- & run-20260607-193053, run-20260607-221822 \\
\addlinespace
H2 mechanism: UDP-pool decomposition & 2 $\times$ $i = 1$ probe runs + 5 s protocol samples & composition: UDP $-$50--58\% vs TCP growth under spread & --- & run-20260610-200013, run-20260610-201131; \texttt{data/conntrack-probe/} \\
\addlinespace
H5: cross-node fraction separates node-local from spreading & 2 $\times$ (8 strategies $\times$ $i = 4$) & lowest-2-of-8 ranks in both batches (joint $\approx$ 0.0013); $\rho = 0.79 \rightarrow 0.25$ & \ref{fig:h5-fraction}, \ref{fig:tradeoff} & run-20260608-070638, run-20260610-202426 \\
\addlinespace
H6: blast radius + recovery trade-off & 2 node-drain batches + 6-strategy gradient & predicted = observed blast every iteration; gradient $\rho = 1.0$ ($n = 6$); recovery contrast at extremes & \ref{fig:h6-trough}, \ref{fig:tradeoff} & run-20260608-194827, run-20260608-205229, run-20260610-172430 \\
\bottomrule
\end{tabular}
\end{center}

\section{What the gate retracted}\label{sec:retractions}

Two of this study's most striking numbers appear nowhere among its claims,
and that is a result in itself. The dirty H4 pilot read as ``co-location $\sim$3$\times$
better at the user layer''; the clean, doctor-gated replication collapsed the
user-facing ratios to $\sim$1.05--1.40$\times$ with no dependency specificity, and the
reading was retracted (\S\ref{sec:h4}). H5's batch-1 continuous correlation (Spearman
$\rho = 0.79$) collapsed to $\rho = 0.25$ n.s.\ in batch 2, and only the two-regime
separation is claimed (\S\ref{sec:h5}). In both cases the provenance-and-replication
gate (\S\ref{sec:campaign}) did exactly what it was designed to do: it caught the most
quotable numbers in the study and removed them. A method that retracts its
own most striking results is the credibility argument for contribution 3 ---
the campaign protocol's discipline is demonstrated here, not merely stated
(\S\ref{sec:not-claimed}, \S\ref{sec:conclusion}).
